\section{Contribuciones del trabajo}

La contribución del trabajo es un procedimiento para aprovechar las respuestas de un dataset en dos etapas de adaptación. Una parte proporciona las demostraciones del entrenamiento supervisado, mientras que otra permite extraer criterios para evaluar las respuestas generadas durante GRPO. La propuesta combina comprobaciones mediante código con un juez semántico y añade ayudas que se retiran durante el entrenamiento, junto con una recompensa que favorece la recuperación de conceptos sin asistencia.

La propuesta se implementa en el dominio clínico y se examina mediante una comparación entre el modelo adaptado con SFT y su posterior optimización con GRPO, utilizando HealthBench como benchmark externo. Este análisis permite identificar cambios en las distintas dimensiones de las respuestas y revisar ejemplos de mejoras y regresiones, considerando las limitaciones del juez local y de la señal utilizada durante el entrenamiento. Además, se conservan configuraciones y registros que permiten inspeccionar las etapas del proceso, mientras que el capítulo~\ref{ch:contribuciones} desarrolla los aportes metodológicos y experimentales junto con sus alcances.
